\SourceLeaf{155}
wenig vergrössert werden müssen, so dass nach längerer Zeit der Spielraum durch 20—34 bezeichnet werden muss. Das ist alles, was die Wahrscheinlichkeitsrechnung lehren kann, \textit{so lange man die Gesellschaft gleichsam nur massenweise betrachtet.} Vielleicht meint mancher, das sei wenig! Ich dächte doch nicht. Es ist sehr wichtig, dass man das Maass der Erwartungen, die man zu haben befugt ist, kennt, und sich nicht einer unbegründeten Sicherheit überlässt. Erwägen wir die beiden Extreme. Es ist \textit{möglich}, dass nach 25 Jahren die Pensionenzahl noch nicht die Zahl 18 überschritten hat, und sich auch noch eine geraume Zeit länger auf dieser oder einer sehr wenig grössern Höhe erhält. Die Pensionenzahl kann auch vorher von ihrer jetzigen Höhe (19) erst noch herabsteigen, und so von jetzt bis 1870 durchschnittlich nicht unbeträchtlich unter 18 sein. Auf diese Weise kann die Kasse Kräfte sammeln, mit denen sie, wenn erst nach langer Zeit das Blatt sich wendet, auch den extremen Fällen der andern Seite die Spitze bieten kann. Das schlimmste wäre eigentlich, wenn diese günstigen Voraussetzungen gar sich noch weiter realisirten; ich meine, wenn in dieser Zwischenzeit die Pensionenzahl erst noch einmal unter 16 herabginge, und so die Kasse in Folge der bestehenden Satzung noch mit neuer Pensionserhöhung von 10, 20, vielleicht gar von 30 Rthl. belastet würde. — Würde, umgekehrt, was aber eben so wohl möglich ist, die Pensionenzahl schon nach 25 Jahren auf oder nahe auf die Höhe von 32 gestiegen sein, so bedarf es nur eines rohen Überschlages, um sich zu überzeugen, dass abgesehen von ganz ausserordentlichen Zuflüssen, die Kasse solchen Anforderungen nicht gewachsen sein, sondern, vermuthlich schon früher, zur Erklärung ihres Unvermögens genöthigt sein wird.

Dass die Überschläge sich etwas günstiger gestalten, wenn man anstatt des Verhältnisses \(7:4\) das von \(2:1\) annimmt, versteht sich von selbst. Man würde dann nach 25 Jahren den Spielraum von 16—28 Pensionen und nach noch längerer Zeit den von 18—30 zu erwarten haben.

Endlich muss ich noch bemerken, dass hiebei stillschweigend vorausgesetzt ist, dass der Umfang der Genossenschaft auf seiner gegenwärtigen Höhe fortan bleibt. Nimmt er noch weiter zu, so muss man sich auf verhältnissmässig noch grössere Zahlen gefasst machen; vermindert er sich hingegen wieder, so wird auch von obigen Resultaten einiger Abzug gemacht werden dürfen. Nach der Natur der Sache aber ist es wohl wenig wahrscheinlich, dass eine \textit{bedeutende} Verminderung anders als nach Verlauf längerer Zeit eintreten könne.

\Needspace{5\baselineskip}\pdfbookmark[2]{Zweite Abtheilung}{witwen-II-2}
\begin{center}\textit{Zweite Abtheilung.}\end{center}
\noindent\hspace*{\parindent}
Da ich die zu der Aufstellung der Bilanz der Witwenkasse einzuschlagenden Wege bereits in meinem Votum vom 8. Januar d. J. umständlich beschrieben, und bei der Ausführung der Arbeit zu Abänderungen keine Veranlassung gefunden habe, so kann ich mich nur auf jenes beziehen, und will hier nur bemerken, dass das Wesen der Methode in der Ermittelung des gegenwärtigen Geldwerths der Obliegenheiten der Kasse besteht, nach den drei Rubriken
\begin{quote}
Obliegenheiten gegen die jetzigen 19 Witwen,\par
Obliegenheiten gegen die Witwen (und Waisen) der jetzigen 51 Theilnehmer,\par
Obliegenheiten gegen die Witwen und Waisen der künftig beitretenden,
\end{quote}
\noindent wobei für die zweite und dritte Rubrik der gegenwärtige Geldwerth der Gegenleistungen durch die Beiträge, in Abzug zu bringen ist. Ich schicke hier zuvörderst einige allgemeine Erläuterungen voraus.

Als Epoche, auf welche alle künftigen Leistungen durch Discontirung bezogen werden, ist der 1. October 1845 gewählt, und vorausgesetzt, dass die Witwen und Interessenten ihre bis dahin fällig gewordenen Pensionen und Beiträge schon resp. empfangen und geleistet haben.
